% \documentclass[aspectratio=169,notes]{beamer}
\documentclass[aspectratio=169]{beamer}
\usetheme[faculty=phil]{fibeamer}
\usepackage{polyglossia}
\setmainlanguage{english} %% main locale instead of `english`, you
%% can typeset the presentation in either Czech or Slovak,
%% respectively.
\setotherlanguages{russian} %% The additional keys allow
%%
%%   \begin{otherlanguage}{czech}   ... \end{otherlanguage}
%%   \begin{otherlanguage}{slovak}  ... \end{otherlanguage}
%%
%% These macros specify information about the presentation
\title[Artificial Intelligence Software (AIS)]{Artificial Intelligence Software, Project 2} %% that will be typeset on the
\subtitle{Inference in the Loop: AI-Driven Robot Control with ROS2
\\ \  \\ \  } %% title page.
\author{Oleg Bulichev}
%% These additional packages are used within the document:
\usepackage{ragged2e}  % `\justifying` text
\usepackage{booktabs}  % Tables
\usepackage{tabularx}
\usepackage{tikz}      % Diagrams
\usetikzlibrary{calc, shapes, backgrounds}
\usetikzlibrary{decorations.pathreplacing,calligraphy,calc,graphs}
\usepackage{amsmath, amssymb}
\usepackage{url}       % `\url`s
\usepackage{listings}  % Code listings
% \usepackage{subfigure}
\usepackage{floatrow}
\usepackage{subcaption}
\usepackage{mathtools}
\usepackage{todonotes}
\usepackage{fontspec}
\usepackage{multicol}
\usepackage{pdfpages}
\usepackage{wrapfig}
\usepackage{qrcode}
\usepackage{animate}
\usepackage{booktabs}
\usepackage{multirow}
% \usepackage{graphicx}
\usepackage{colortbl}

\graphicspath{{resources/}}
\frenchspacing

\setbeamertemplate{caption}[numbered]
\usetikzlibrary{graphs}

% \usepackage[backend=biber,style=ieee,autocite=footnote]{biblatex}
% \addbibresource{biblio.bib}
% \DefineBibliographyStrings{english}{%
%   bibliography = {References},}

\newcommand{\oleg}[2][] {\todo[color=red, #1] {OLEG:\\ #2}}
\newcommand{\fbckg}[1]{\usebackgroundtemplate{\includegraphics[width=\paperwidth]{#1}}}%frame background

\usepackage[framemethod=TikZ]{mdframed}
\newcommand{\dbox}[1]{
\begin{mdframed}[roundcorner=3pt, backgroundcolor=yellow, linewidth=0]
\vspace{1mm}
{#1}
\vspace{1mm}
\end{mdframed}
}

\begin{document}
\setlength{\abovedisplayskip}{0pt}
\setlength{\belowdisplayskip}{0pt}
\setlength{\abovedisplayshortskip}{0pt}
\setlength{\belowdisplayshortskip}{0pt}

\fbckg{fibeamer/figs/title_page.png}
\frame[c]{\setcounter{framenumber}{0}
    \usebeamerfont{title}%
    \usebeamercolor[fg]{title}%
    \begin{minipage}[b][6.5\baselineskip][b]{\textwidth}%
        \textcolor{black}{\raggedright\inserttitle}
    \end{minipage}
    % \vskip-1.5\baselineskip

    \usebeamerfont{subtitle}%
    \usebeamercolor[fg]{framesubtitle}%
    \begin{minipage}[b][3\baselineskip][b]{\textwidth}
        \raggedright%
        \insertsubtitle%
    \end{minipage}
    \vskip.25\baselineskip
}
%   \frame[c]{\maketitle}

\fbckg{fibeamer/figs/common.png}

\note{\scriptsize \begin{itemize}
        \item \
    \end{itemize}}

\begin{frame}[t]{Short Description}
    \vspace*{-0.5cm}
The goal of this project is to build an \textbf{AI-driven robot control system} by integrating a trained ML model with a ROS2 pipeline using \textbf{ONNX} as the inference format.

The full pipeline:
\begin{multicols}{2}
\begin{enumerate}
    \item Define an AI task and a robot action mapping
    \item Train or fine-tune a model
    \item Export the model to \textbf{ONNX}
    \item Deploy inference inside a \textbf{ROS2 node}
    \item Control a robot via ROS topics based on model output.
    \item Monitor live inference in a \textbf{Gradio} panel
    \item Demonstrate in \textbf{Gazebo} or on a real robot
\end{enumerate}
Students are free to choose any modalities like: CV, audio, NLP, etc.
\end{multicols}
Reusing a model from \textbf{Project~1} is encouraged but not required.
\end{frame}


\begin{frame}[t]{Example Project}
\framesubtitle{Gesture-Controlled Differential-Drive Robot}
\vspace{-0.3cm}
A YOLO classification model recognize \textbf{hand gestures} from a webcam.
Each gesture class is mapped to a robot motion command:
        
\textbf{Gesture classes (example):}
        \begin{itemize}
            \item \texttt{forward} $\to$ move forward
            \item \texttt{backward} $\to$ move backward
            \item \texttt{turn\_left} $\to$ rotate left
            \item \texttt{turn\_right} $\to$ rotate right
            \item \texttt{stop} $\to$ zero velocity
        \end{itemize}
\vspace{0.2cm}
The robot publishes velocity commands to \texttt{/cmd\_vel}.
A \textbf{Gradio} panel displays the live predicted class and confidence score.
\end{frame}


\begin{frame}[t]{Step 1}
\vspace{-0.3cm}
\textbf{Short description}: Define the AI task your system will solve and specify how the model output maps to robot actions.

\textbf{Artifacts}: Report section describing the task design.

\textbf{Requirements}
\footnotesize
\begin{itemize}
    \item Describe the input modality (camera, microphone, etc.)
    \item Specify the model output: class labels, bounding boxes, regression values, etc.
    \item Define the complete \textbf{class-to-action mapping}\\
          (e.g.\ class \texttt{forward} $\to$ \texttt{linear.x = 0.3})
    \item Justify why the chosen task is suitable for robot control
    \item The system must support at least \textbf{3 distinct robot actions}
\end{itemize}
\end{frame}

\begin{frame}[t]{Step 2}
\vspace{-0.3cm}
\textbf{Short description}: Obtain a trained AI model that produces the outputs defined in Step~1 and is exportable to ONNX.

\textbf{Artifacts}: Model checkpoint (or downloaded weights), evaluation results, report section.

\textbf{Requirements}
\footnotesize
\begin{itemize}
    \item Use any ONNX-exportable framework (Ultralytics YOLO, PyTorch, HuggingFace, etc.)
    \item \textbf{Fine-tuning is not required} --- you may use a pre-trained model directly if it already solves your task
    \item If the available pre-trained models do not fit, fine-tune on a suitable dataset
    \item Evaluate the chosen model with appropriate metrics (accuracy, mAP, F1, etc.)
    \item Describe the model selection or training process in the report
    \item ClearML experiment tracking is \textbf{recommended} if training is performed
\end{itemize}
\end{frame}

\begin{frame}[t]{Step 3}
\vspace{-0.3cm}
\textbf{Short description}: Export the trained model to ONNX format and validate that inference produces results consistent with the original framework.

\textbf{Artifacts}: \texttt{.onnx} model file, validation script, report section with comparison results.

\textbf{Requirements}
\footnotesize
\begin{itemize}
    \item Export the model to ONNX using the framework's export API
    \item Document: input/output tensor names and shapes, opset version, dynamic axes if used
    \item Run a validation script: pass the same input through both the original model and \texttt{onnxruntime}
    \item Verify outputs match within floating-point tolerance
    \item Measure and report inference latency for the ONNX model on your hardware
\end{itemize}
\end{frame}

\begin{frame}[t]{Step 4}
\vspace{-0.4cm}
\textbf{Short description}: Design and document the ROS2 system architecture for your project.

\textbf{Artifacts}: Architecture description in the report with all node and topic names defined.

\textbf{Requirements}
\footnotesize
\begin{itemize}
    \item You choose how inference integrates with ROS2:
    \begin{itemize}\footnotesize
        \item \textbf{Outside ROS2} --- a standalone script reads the sensor directly and publishes the result to a ROS2 topic
        \item \textbf{Inside ROS2} --- a dedicated ROS2 node subscribes to a sensor topic, runs ONNX, and publishes the result
    \end{itemize}
Either way, the ROS2 side must include at minimum: a \textbf{result topic} carrying the model output (e.g.\ \texttt{/model\_result}); a \textbf{controller node} that subscribes to the result and commands the robot
    \item List all ROS2 topics: name, message type, publisher, subscriber
    \item Justify your integration choice in the report
    \item Describe the complete data flow: sensor $\to$ inference $\to$ ROS2 topic $\to$ robot actuator
\end{itemize}
\end{frame}

\begin{frame}[t]{Step 5}
\vspace{-0.3cm}
\textbf{Short description}: Implement the inference component that runs the ONNX model and publishes the result to a ROS2 topic.

\textbf{Artifacts}: Inference script or ROS2 node source code.

\textbf{Requirements}
\footnotesize
\begin{itemize}
    \item Implement according to the architecture chosen in Step~4:
    \begin{itemize}\footnotesize
        \item \textbf{Outside ROS2}: a Python script reads sensor data directly (e.g.\ via OpenCV, PyAudio), runs ONNX inference, and publishes to \texttt{/model\_result} via \texttt{rclpy}
        \item \textbf{Inside ROS2}: a ROS2 node subscribes to a sensor topic, runs ONNX inference, and publishes to \texttt{/model\_result}
    \end{itemize}
    \item Load the ONNX model using \texttt{onnxruntime} at startup
    \item Preprocess input data to match the model's expected input shape
    \item Published message must contain: predicted class name and confidence score
\end{itemize}
\end{frame}

\begin{frame}[t]{Step 6}
\vspace{-0.3cm}
\textbf{Short description}: Implement a ROS2 controller node that subscribes to the model result and commands the robot accordingly.

\textbf{Artifacts}: ROS2 package containing the controller node source code.

\textbf{Requirements}
\footnotesize
\begin{itemize}
    \item Subscribe to \texttt{/model\_result} published by the inference script
    \item Map each predicted class or value to an appropriate ROS2 command message.\\
          The command type depends on the robot:
    \begin{itemize}\footnotesize
        \item mobile robot: \texttt{geometry\_msgs/Twist} on \texttt{/cmd\_vel}
        \item manipulator: \texttt{trajectory\_msgs/JointTrajectory}
        \item custom actuator: any suitable message type
    \end{itemize}
    \item All action-to-command mappings must be explicitly defined and documented
    \item Include a configurable \textbf{confidence threshold} --- send a safe/neutral command if confidence is below it
\end{itemize}
\end{frame}

\begin{frame}[t]{Step 7}
\vspace{-0.3cm}
\textbf{Short description}: Build a Gradio application that subscribes to the inference result topic and displays a live monitoring panel.

\textbf{Artifacts}: Python source code for the Gradio monitoring app.

\textbf{Requirements}
\footnotesize
\begin{itemize}
    \item The panel must display in real time:
    \begin{itemize}\footnotesize
        \item current predicted class label
        \item confidence score
        \item optionally: the latest input frame
    \end{itemize}
    \item Subscribe to \texttt{/model\_result} via \texttt{rclpy} or a shared queue
    \item Use appropriate Gradio components (\texttt{gr.Label}, \texttt{gr.Image}, etc.)
    \item The panel must update continuously without manual refresh
    \item The app runs as a \textbf{separate process} alongside the ROS2 nodes
\end{itemize}
\end{frame}

\begin{frame}[t]{Step 8}
\vspace{-0.3cm}
\textbf{Short description}: Demonstrate the complete integrated system with the robot responding to live model inference while the Gradio panel monitors inference output.

\textbf{Artifacts}: Recorded demo video; report section describing the demo setup.

\textbf{Requirements}
\footnotesize
\begin{itemize}
    \item All components must run simultaneously and be visible in the recording:
    \begin{itemize}\footnotesize
        \item live sensor input
        \item robot reacting in \textbf{Gazebo simulation} or on a \textbf{real robot}
        \item \textbf{Gradio panel} updating in real time
    \end{itemize}
    \item The video must demonstrate at least \textbf{3 distinct robot actions} triggered by model output
    \item Describe the simulation world or real robot setup in the report
\end{itemize}
\end{frame}


\begin{frame}[t]{Submission}
You must submit a \textbf{report} and a working code repository covering the full pipeline.
\vspace*{-0.3cm}
\begin{multicols}{2}
The report must include:
\begin{itemize}
    \item task design and action mapping (Step~1)
    \item model training process and results (Step~2)
    \item ONNX export settings and validation (Step~3)
    \item ROS2 architecture description (Step~4)
    \item description of all implemented nodes (Steps~5--6)
    \item Gradio panel description and screenshots (Step~7)
    \item demo setup and observations (Step~8)
    \item discussion of difficulties and limitations
\end{itemize}
Required artifacts:
\begin{itemize}
    \item Git repository with all source code
    \item trained \texttt{.onnx} model file
    \item demo video (robot + Gradio panel)
    \item PDF report
\end{itemize}
\end{multicols}
\end{frame}




\fbckg{fibeamer/figs/last_page.png}
\frame[plain]{}

\end{document}